We studied status updates from a UAV whose sensor is its own classifier, tracking the GS's AoII through a belief that each label updates according to its confidence. This belief gives closed-form one-step values for the three actions and a drift-plus-penalty policy with confidence thresholds. On flood-mapping traces under chip-level cross-validation, the policy outperforms adaptations of three recent designs at every budget up to $5$\,mJ per slot. Against an AoII-threshold policy it helps when energy is scarce or the on-board classifier is weak, because then the confidence decides which disagreements are worth reporting. With an accurate classifier and ample energy, reporting every disagreement works as well, and a belief that knows only the precision of each label value matches one that uses the per-frame confidence.

On-board errors on overlapping footprints are correlated in time, and our belief handles this only through a tempered likelihood. A belief that models the correlation is the most promising way to improve the policy at large budgets. The one-step rule also does not anticipate predictable outages such as range limits or satellite passes, and the structure of the optimal POMDP policy and a performance bound for a renewal version still need proofs. Finally, Sentinel-2 imagery, generated separately from the UAV position, stands in for UAV frames. A spatial model of the area, with revisits and a route chosen by the UAV, is future work. Repeating the study with UAV datasets such as FLAME~\cite{shamsoshoara2021flame} and AIDER~\cite{kyrkou2020aider}, with energy measured on a microcontroller, is the next step.
